\clearpage
\section{Hardware Software Contracts}
\label{app:hwsfcontract}

We succinctly repeat the definitions and framework for Guarnieri et. al..
Full details can be found in~\cite{guarnieri2021hardware}.
% Let's $p$ be a program.
% For a given CPU, we define $\sigma$, the architectural state composed of a memory and a set of registers.
% Details on how to model the architectural semantics i.e. the transition between architectural states can be found in~\cite{guarnieri2021hardware}.
% Hardware state is defined by combining a architectural state $\sigma$ and a microarchitectural state $\mu$, and a hardware semantic is used to express transition between hardware states and their successors.
\\

\noindent\textbf{Hardware-Software Contracts}
Contracts define public or exposed states with respect to \emph{programs}.
For a given program $p$ and an initial architectural state $\sigma_0$, $\contract[]{p}{}(\sigma_0)$ is the traces of architectural state when the program is run given the corresponding semantic.
The traces of a contract $\contract[]{p}{})$ capture which architectural state is public and at risk of being exposed by the software when run.
\\

\subsection{Hardware Semantics}
The \emph{adversary model} is modeled as projections of parts of the microarchitectural state.
Given a program $p$, and an architectural state $\sigma$, $\uarch[]{p}{}(\sigma)$ denotes the trace of hardware observations produced in the run of the program on the given microarchitecture.
The hardware semantic of a program $p$ is denoted as $\sigma$, $\uarch[]{p}{}$.
\\

\noindent\textbf{Definition 1} ($\uarch[]{\cdot}{} \vdash \contract[]{\cdot}{}$)
A hardware semantics $\uarch[]{\cdot}{}$ \emph{satisfies a contract} $\contract[]{\cdot}{}$ if, for all programs $p$ and all initial architectural states $\sigma$, $\sigma'$, if $\contract[]{p}{}(\sigma) = \contract[]{p}{}(\sigma')$ then $\uarch[]{p}{}(\sigma) = \uarch[]{p}{}(\sigma')$.

\subsection{Building Blocks for Contracts}
Contracts for secure speculation are defined two fold.
They are composed of a \emph{leakage model} and an \emph{execution model}.
The leakage model describes what is observable by the attacker (e.g. execution trace and trace of memory accesses).
The execution model, on the other hand, models the speculative behavior of the underlying semantic.
\\

\noindent\textbf{Existing Leakage Models}
Several leakage models have been defined and used in previous work.
The \emph{constant-time} model (\sw{ct} for short) leaks the trace of the \texttt{pc} (i.e. the control flow) and the trace all memory accesses (but not their values).
The \emph{architectural} or \emph{sanboxing} model (\sw{arch}) extends the \sw{ct} model to include the values of elements loaded from memory.
\sw{arch} effectively exposes the value of all registers during execution.
The \emph{memory} model (\sw{mem}) only exposes the trace of memory accesses but not the control flow of the program and is sometimes used to model an attacker exploiting cache side-channels.
The more states are exposed to an attacker by a leakage model, the \emph{stronger} it is considered to be (i.e. challenging to defend by software-only measures).
Hence \sw{arch} being the strongest leakage model followed by \sw{ct} and finally \sw{mem}.
\\

\noindent\textbf{Existing Execution Models}
The weaker execution model is \sw{seq}, which represents non-speculative, or sequential execution (one instruction at a time and in-order).
We also have the strongest execution model \sw{spec} which models speculative execution with the following predictors: the pattern history table (PHT), the branch target buffer (BTB) and the return stack buffer (RSB).
\\

\subsection{SW side: Non-Interference Properties}

Work by Cauligi et. al.~\cite{cauligi2022sok} defines direct and relative non-interference properties.
We repeat these definitions here as a reminder.
\\

\noindent\textbf{Direct Non-Interference.}
Program $p$ satisfies \emph{direct non interference} with respect to a given contract $\contract[]{\cdot}{}$ and policy $\pi$ if, for all pair of $\pi$-equivalent initial states $\sigma$ and $\sigma'$, executing $p$ with each initial trace produces the same trace. That is, $p \vdash NI(\pi, \contract[]{\cdot}{})$ is defined as
$$ \forall \sigma, \sigma': \sigma \simeq_{\pi} \sigma' \Rightarrow \contract[]{p}{}(\sigma) = \contract[]{p}{}(\sigma')$$
We elide $\pi$ for brevity and write $p \vdash NI(\contract[]{\cdot}{})$.
\\

\noindent\textbf{Removing Security Policy $\pi$.}
In Cauligi et. al.~\cite{cauligi2022sok}'s definitions of relative non-interference properties all contain a security policy $\pi$, which is used to qualify which part of the initial states $\sigma$ and $\sigma'$ are public or private.
This is redundant for the definition or \emph{relative} non-interference as what is public is already defined by what is leaked in the first contract.
As a result, we elide $\pi$ for brevity.
\\

\noindent\textbf{Relative Non-Interference.}
Program $p$ satisfies \emph{relative non interference} from contract $\contract[a]{\cdot}{seq}$ to $\contract[a]{\cdot}{\beta}$ if: for all pair of initial state $\sigma$ and $\sigma'$, if executing $p$ under $\contract[a]{\cdot}{seq}$ produces equal traces, then executing $p$ under $\contract[a]{\cdot}{\beta}$ produces equal traces.
That is, $p \vdash NI(\contract[a]{\cdot}{seq} \Rightarrow \contract[a]{\cdot}{\beta})$ is defined as
\begin{align*}
\forall \sigma, \sigma': \quad
&\contract[a]{\cdot}{seq}(\sigma) = \contract[a]{\cdot}{seq}(\sigma') \\[1ex]
&\Rightarrow \contract[a]{\cdot}{\beta}(\sigma) = \contract[a]{\cdot}{\beta}(\sigma')
\end{align*}

\section{Burst Mode Security Analysis}
\label{app:proofburstmode}
Given that $\uarch[burst]{\cdot}{} \vdash \contract[shm]{\cdot}{stl}$.
Given that our static program analysis tool \hw{sta} returns \texttt{true} on a program $p$ if $p \vdash NI(\contract[shm]{\cdot}{seq} \Rightarrow \contract[shm]{\cdot}{stl})$, and \texttt{false} otherwise.
Let us prove $\uarch[burst]{p}{sta} \vdash \sw{(}\contract[shm]{\cdot}{seq}$.
\\

Let us take $p$ a program, $\sigma$ and $\sigma'$ some initial state so that
$(\contract[shm]{p}{seq}(\sigma) = \contract[shm]{p}{seq}(\sigma'))$
We have two cases:
\\

\noindent
If $sta(p) = \texttt{true}$, $p \vdash NI(\contract[shm]{\cdot}{seq} \Rightarrow \contract[shm]{\cdot}{stl})$ and $\uarch[burst]{p}{} = \uarch[burst]{p}{sta}$ so:
\begin{align*}
    (\contract[shm]{p}{seq}(\sigma) &= \contract[shm]{p}{seq}(\sigma')) \\[1ex]
    \Rightarrow &(\contract[shm]{p}{stl}(\sigma) = \contract[shm]{p}{stl}(\sigma')) \\[1ex]
    \Rightarrow &  \uarch[burst]{p}{}(\sigma) = \uarch[burst]{p}{}(\sigma') \\[1ex]
    \Rightarrow &  \uarch[burst]{p}{sta}(\sigma) = \uarch[burst]{p}{sta}(\sigma')
\end{align*}

\bigskip
\noindent
If $sta(p) = \texttt{false}$, $\uarch[burst]{p}{sta} = \varnothing$ and:
$$\uarch[burst]{p}{sta}(\sigma) = \uarch[burst]{p}{sta}(\sigma') \text{ is always true.}$$

\bigskip
\noindent
As a result, under our threat model, $\uarch[burst]{p}{sta} \vdash \contract[shm]{\cdot}{seq}$.